\subsection{DAQP with side-specific elastic weights}\label{sec:daqp}
DAQP uses a dual active-set strategy for convex quadratic programs and exposes
side-specific soft weights through its public interface~\cite{daqp,goldfarb}.
This capability is important to the comparison: asymmetry alone cannot be
claimed as a new modeling feature of an original-row projection method. A fair
adapter must reproduce the requested finite objective rather than select the
backend's default softness convention and call the result equivalent.

For a compatible elastic side with departure $t\geq0$, the requested price is
\begin{equation}\label{eq:daqp-cost}
 \varphi(t)=a t+\tfrac12 b t^2,\qquad a\geq0,\quad b>0.
\end{equation}
The adapter maps the quadratic convention reciprocally where the public
backend defines a softness parameter through its inverse curvature. The precise
mapping is applied per side, with lower and upper coefficients kept separate.
It is validated by evaluating \eqref{eq:daqp-cost} in original coordinates after
the solve, rather than assuming that similarly named interface fields have the
same units or normalization. Dividing by a zero quadratic price is not a valid
way to request a pure linear objective.

An active set identifies constraints whose multipliers may be nonzero. For a
positive-definite original Hessian, eliminating the primal vector yields a dual
quadratic structure involving
\begin{equation}\label{eq:daqp-dual}
 x=-P^{-1}(q+A^\top y),\qquad
 A P^{-1}A^\top,
\end{equation}
with the finite-side and elastic conventions handled by the backend. This
identity explains why positive-definite curvature matters to the direct path;
it is not a prescription to form dense inverses explicitly. The tested adapter
uses the library interface and its maintained factors. It does not repair a
semidefinite scientific model by adding an undisclosed Hessian shift.

The ordinary comparison admits DAQP when the base Hessian and finite soft-side
quadratic coefficients meet the implemented positive-definite elastic domain.
Pure L1 and grouped/Huber objectives are incompatible with that particular
mapping. The coupled study has separately registered direct epigraph and
proximal DAQP policies, including their stated linear-programming cases. Those
policies are not silently substituted into the earlier experiment. Distinct
protocols make their different acceptance and termination conditions inspectable.

Algorithm~\ref{alg:daqp} describes the direct side-weight adapter. The conversion
has explicit size bounds of 2048 variables and 8192 finite sides or rows where
applicable. These are implementation resource limits, not mathematical
restrictions of convex quadratic programming. A refusal at this boundary is
retained as an unavailable operation; it does not justify reporting a successful
answer to a smaller or differently regularized model. Inputs that pass the size
check still require semantic compatibility and original-space evaluation.

\begin{algorithm}[htbp]
\caption{Map compatible side-specific elastic prices into DAQP, retain its active workspace for allowed vector updates and recover original-space objective and multipliers.}\label{alg:daqp}
\begin{algorithmic}[1]
\Require Positive-definite admitted QP and finite-side elastic coefficients
\State Check finite dimensions, hard mask, side pattern and strictly positive quadratic prices
\State Convert each lower and upper softness convention with its physical scale
\State Prepare DAQP; retain factor and active state for compatible vector updates
\State Solve; interpret the backend's own status convention
\State Recover original dual signs and independently evaluate all original rows
\State Return accepted only if the external accuracy requirements also hold
\end{algorithmic}
\end{algorithm}

Repeated-operation fairness requires using DAQP's documented update operation.
The prepared workspace preserves its active state and updates only the allowed
linear term and bounds. It does not resend a complete sense vector on every
frame when the public API already permits a cheaper vector update. Conversely,
changing penalty structure invalidates that reuse. The measured timing includes
the selected preparation/update boundary consistently, so the comparison is not
an active-set warm start against a deliberately rebuilt reference.

Backend status interpretation is method-specific. DAQP's status indicating a
soft optimum cannot be read using OSQP's integer convention, where a numerically
equal value may denote an inaccurate result. The adapter translates the status,
then recomputes objective discrepancy and maximum raw hard violation. In the
coupled study, candidate vectors also face the separate primal and dual checks.
An active-set result that is fast but solves the wrong softness scaling is
therefore rejected; a correct result outside the supported direct domain is not
manufactured by altering the objective. These controls preserve the scientific
comparison even when different backends make different engineering tradeoffs.
The separately registered proximal DAQP policies address additional epigraph
regimes through a proximal-point interpretation rather than a permanent
scientific Hessian shift~\cite{daqpprox,rockafellarppa}. Their original objective
is checked after numerical recovery. This distinction matters for a zero-
Hessian linear-price case: a positive quadratic term introduced only to
make a direct solver callable would change the requested model. The protocol
instead admits only its explicitly specified policies and retains unsupported
or unsuccessful outcomes. These later grouped policies do not retroactively
expand the direct adapter's compatibility in the original separable timing
study, whose method settings and excluded cells remain fixed.
